\section{Event account against spectral gap}
\label{sec:h-merge}

Each merge has the event-local value
\[
R=A_{\mathrm{coef}}+\frac{|d|^2}{2}+\Re(\bar d\, W)
=-(\Delta E_\psi+\Delta E_G),
\]
and the reverse split carries \(-R\).
Along a sequence, \(\sum_i R_i\) equals the change in known energy, because
\(E_\psi\) is conserved between events.
Radiation that leaves the region is part of that final known energy.
The gap \(m=-8-E_0\) of Proposition~\ref{prop:clique} depends only on the
final graph and the final eigenmode.
Two sequences that arrive at the same \((G,\psi)\) share \(m\) and need not
share \(\sum R\).
Norm change of the same event is \(2B\), which is Section~\ref{sec:bond}
and is not a firing condition.
